\scanpage{63}
\mspage{69}

c) La formation de $f_*F$ \inserted{($F$ \deleted{\unsure{$\in\mathrm{Ob}$}} de $\unsure{\mathbb D|\underline D}(X)$)} commute à tout changement de base $Y'\longrightarrow Y$, et ça reste vrai \deleted{\unsure{pour}} après \deleted{tout} chgt de base -- i.e. si on a $Y''\longrightarrow Y'\longrightarrow Y$ d'où

\diagram{\begin{tikzcd}[ampersand replacement=\&,row sep=tiny,column sep=small]
X'' \arrow[r] \arrow[d,"f''"'] \& X' \arrow[r] \arrow[d,"f'"] \& X \arrow[d,"f"] \\
Y'' \arrow[r,"p'"'] \& Y' \arrow[r,"p"'] \& Y
\end{tikzcd}}

, ``le th. de chgt de base'' est vrai pour $f'_*$ et $p''^*$.

% slanted
\margin{c$'$) La formation de $f_*(F)$ ($F\in\mathrm{Ob}\,\mathbb D(X)$) commute au passage aux fibres, et ça reste vrai après \deleted{tout} chgt. de base $p:Y'\longrightarrow Y$. c$''$) Comme c$'$), avec $Y'\longrightarrow Y$ \unsure{une injection|un morphisme|une ingression} de localisation $Y_{/y}\longrightarrow Y$. d) Si $Y''\xrightarrow{p}Y'\longrightarrow Y$, avec $p$ \unsure{coasph.|coasphér.}, alors $q:X''=X\times_Y Y''\longrightarrow X'=X\times_Y Y'$ \unsure{coasphérique} (\deleted{\unsure{donc}} \unsure{Utilisation|Utiliser} \unsure{du th. 3}). d$'$) Comme d), avec $Y''=e$ (donc) $Y'$ cat. avec objet final $e$ \unsure{$/p$}). \unsure{d$'$)|d$''$)} Comme d), avec $Y''=\unsure{\{i\}|\{1\}|\{e\}}$ et $Y'=\Delta^1$ \unsure{$(p=\{0,1\})$}}

\uline{Dém.} Pour a) $\Longleftrightarrow$ b) cf prop. page 53. Montrons que \deleted{Ces conditions} \deleted{équivalentes} ces conditions impliquent c). Comme elles \unsure{sont} stables par changement de base (cf dém. \deleted{prop.} \inserted{immédiate} \unsure{p.} 55), donc \inserted{par Der 2} on est ramené à prouver \deleted{\unsure{l'iso}} \deleted{que} \deleted{que pour} \unsure{que}

% slanted
\margin{Donc on utilise \uline{Der 5 (d)}}

\[ H^\bullet_{\mathbb D}(X_{/y},F)\mathrel{\overset{\sim}{\unsure{\longleftrightarrow|\longrightarrow}}}H^\bullet_{\mathbb D}(Y_{/y},F)\qquad(F\in\mathrm{Ob}\,\mathbb D(Y)), \]

ce qui résulte du \uline{th. 1}*. (Cf dém. du th. p. \unsure{56|16}, qui marche \unsure{avec ces|avec des} hyp. plus générales.)

On a $c\underset{\vphantom g}{\overset{\text{par Der (2)}}{\Longleftrightarrow}}c'\Longrightarrow c''$

\deleted{Prouvons $c''\Longrightarrow a$. L'hyp. de chgt de base pour $\{y_0\}\longrightarrow Y$, plus \deleted{\unsure{Der}} calcul de la fibre de \deleted{\unsure{$R$}} $f_*(F)$ en $y$ (\uline{Der 5 (g)}), donne que}

\[ \deleted{H^\bullet_{\mathbb D}(X_{/y_0},\deleted{\unsure{F}}G)\longrightarrow H^\bullet_{\mathbb D}(X_{y_0},\varphi^*G)} \]

\deleted{est iso, chaque fois que $G$ est un objet de $\mathbb D(X_{/y_0})$ qui provient d'un $F$ dans $\mathbb D(X)$ par \deleted{\unsure{le foncteur}} \deleted{\unsure{$p$}} $q^*$, où}

\[ \deleted{q:X_{/y_0}\longrightarrow X} \]

\deleted{est le foncteur canonique. Donc on gagne par le th. 1, si on prouve que $q^*$ est ess. surjectif, Il suffit de prouver que pour $G$ sur $X_{/y_0}$, on a}

\[ \deleted{q^*q_*(G)\xrightarrow{\sim}G\qquad\text{(adjonction)}} \]

\deleted{Or pour $(x_0,u_0)\in\mathrm{Ob}(X_{/y_0})$, \inserted{($u_0:f(x_0)\longrightarrow y_0$)} on a, comme $q(x_0,u_0)=x_0$}

\[ \deleted{q^*(q_*(G))_{(x_0,u_0)}=q_*(G)_{x_0}\simeq H^\bullet_{\mathbb D}((X_{/y_0})_{/x_0})} \]
% archive sheet number: 62
